% ------------------------------------------------------------------------------
\section{Actualizar a una versión nueva}
% ------------------------------------------------------------------------------

Cuando se publique una versión nueva:

\begin{enumerate}
	\item Cierre ConcretoLab si está abierto.
	\item Descargue el instalador con el mismo enlace de la sección~\ref{sec:descarga}: siempre entrega la versión
	      más reciente.
	\item Haga doble clic en \ruta{\archivoinstalador} y siga los mismos pasos de la instalación.
	\item El instalador reemplaza la versión anterior; no hace falta desinstalarla antes.
\end{enumerate}

Sus proyectos y copias de seguridad se conservan. Para ver qué versión tiene instalada, abra ConcretoLab y pulse
\menu{Ayuda} \textgreater{} \menu{Acerca de ConcretoLab}.

\begin{consejo}
Antes de actualizar, cree una copia de seguridad desde \menu{Datos y copias} \textgreater{} \boton{Crear copia ahora}.
No es obligatorio, pero es una buena costumbre antes de cualquier cambio importante.
\end{consejo}

% ------------------------------------------------------------------------------
\section{Desinstalar}
% ------------------------------------------------------------------------------

\paso{1}{Abra la lista de aplicaciones instaladas}

\begin{itemize}
	\item \textbf{Windows 11:} \menu{Inicio} \textgreater{} \menu{Configuración} \textgreater{} \menu{Aplicaciones}
	      \textgreater{} \menu{Aplicaciones instaladas}.
	\item \textbf{Windows 10:} \menu{Inicio} \textgreater{} \menu{Configuración} \textgreater{} \menu{Aplicaciones}
	      \textgreater{} \menu{Aplicaciones y características}.
\end{itemize}

\paso{2}{Busque ConcretoLab y desinstálelo}

Escriba \textbf{ConcretoLab} en el buscador de la lista. En Windows~11 pulse los tres puntos \boton{\ldots} que están a la
derecha y elija \menu{Desinstalar}; en Windows~10 haga clic sobre ConcretoLab y pulse \boton{Desinstalar}. Confirme con
\boton{Desinstalar}.

\begin{importante}[Sus proyectos no se borran]
Desinstalar quita el programa y sus accesos directos, pero \textbf{no borra sus proyectos ni sus copias de seguridad}:
quedan en \ruta{C:\textbackslash Users\textbackslash <su usuario>\textbackslash .concretolab}. Si vuelve a instalar
ConcretoLab, los encontrará tal como estaban.

Solo si quiere eliminarlos para siempre, borre esa carpeta a mano después de desinstalar. Antes, exporte o copie a
otro lugar lo que quiera conservar: una carpeta borrada no se puede recuperar desde ConcretoLab.
\end{importante}
